% !TeX root = ../../../geos_mpm_manual.tex
\section{Experimental weak-discontinuity projection}
\label{sec:weak-trace-implementation}
\index{weak discontinuity projection}
\index{weak trace projection}
\index{weak discontinuity}
\index{SurfaceFlag!WeakDiscontinuity}

\noindent\fbox{\begin{minipage}{0.95\linewidth}\small
The weak-discontinuity projection is an \emph{experimental} GEOS-MPM feature.
It is disabled by default and is currently restricted to serial, CPU,
plane-strain calculations with a prescribed particle-carried interface.  The
implementation described here reconstructs a continuous piecewise-linear
interface and weakly ties two contact-group velocity fields with local mortar
projections.  It is not yet a validated replacement for the cohesive-zone
implementation in Section~\ref{sec:cohesive-zone-implementation} or for
ordinary contact in Section~\ref{sec:contact-options}.
\end{minipage}}\par\medskip

\subsection{Purpose and enriched kinematics}
\index{mixed cell}
\index{weak discontinuity!mixed-cell stress error}

A bonded material interface is a weak discontinuity: velocity is continuous,
but its gradient and the material response can differ on the two sides,
\begin{equation}
  \llbracket \mathbf{v} \rrbracket_\Gamma = \mathbf{0},
  \qquad
  \llbracket \nabla \mathbf{v} \rrbracket_\Gamma \ne \mathbf{0}.
  \label{eq:weak-trace-weak-discontinuity}
\end{equation}
Mapping particles with dissimilar constitutive responses to one smooth nodal
velocity field can therefore impose an incorrect mixed-cell strain increment.
The resulting error can be especially severe when the correct response is
series-like in the interface-normal direction.  This is the same motivation
that underlies weak-discontinuity enrichments and the CPDI2 construction
\cite{crook2025cohesive,sadeghirad2013cpdi2}.

The present implementation represents the two material sides with independent
nodal fields,
\begin{equation}
  \mathbf{v}_h^A(\mathbf{x}) = \sum_I N_I(\mathbf{x})\mathbf{v}_I^A,
  \qquad
  \mathbf{v}_h^B(\mathbf{x}) = \sum_I N_I(\mathbf{x})\mathbf{v}_I^B.
  \label{eq:weak-trace-two-fields}
\end{equation}
The fields are created by assigning the materials to distinct prescribed
contact groups.  Their gradients can evolve independently, while a mortar
constraint enforces only the physically required velocity trace on the
reconstructed interface.  No constitutive tangent and no cohesive stiffness
are required.

\subsection{Particle-carried interface data}
\index{contact groups!weak discontinuity projection}
\index{surface reconstruction!weak discontinuity}

Particles adjacent to the bonded interface carry a surface-position vector and
surface normal and use
\begin{equation}
  \texttt{SurfaceFlag::WeakDiscontinuity}=5.
\end{equation}
For particle $p$, the current surface sample is
\begin{equation}
  \mathbf{y}_p = \mathbf{x}_p + \mathbf{s}_p.
  \label{eq:weak-trace-surface-sample}
\end{equation}
The flag supplies geometry without marking the surface as damaged or opened.
For a selected contact-group pair $(A,B)$, normals on the $B$ side are reversed
before fitting.  Consequently, the usual input convention is that each side's
surface normal points outward from its own material; the two coincident
interface normals then become consistently oriented during reconstruction.

The current implementation requires prescribed surface geometry,
\texttt{damageFieldPartitioning=0}, and
\texttt{weakDiscontinuityInsertion=off}.  Automatic material-ID or
stress-residual reconstruction can be added later without changing the mortar
projection itself.

\subsection{Continuous P1 interface reconstruction}
\index{surface reconstruction!P1 marching triangles}
\index{weak discontinuity projection!C0 interface}

At each in-plane background-grid vertex $I$, the solver fits a
signed-distance-like value from the $k$ nearest oriented samples,
\begin{equation}
  \phi_I =
  \frac{\displaystyle\sum_{p\in\mathcal{K}_I}
    w_{Ip}(\mathbf{x}_I-\mathbf{y}_p)\mathbin{\cdot}\widetilde{\mathbf{n}}_p}
       {\displaystyle\sum_{p\in\mathcal{K}_I} w_{Ip}},
  \qquad
  w_{Ip} =
  \frac{1}{\lVert\mathbf{x}_I-\mathbf{y}_p\rVert^2+(h_\Gamma/4)^2},
  \label{eq:weak-trace-nodal-phi}
\end{equation}
where $h_\Gamma=\sqrt{h_xh_y}$ and
$k=\texttt{weakInterfaceTraceSurfaceFitNeighbors}$.  The fitted nodal values
are not stored as an advected material field; they are reconstructed from the
current particle surface samples at every explicit step.

Every Cartesian cell is split along the same $p_{10}$--$p_{01}$ diagonal.  A
P1 scalar is defined on each triangle from its three values of $\phi_I$, and
the zero contour is intersected with the triangle edges.  Vertices on shared
grid or simplex edges are keyed by their common endpoint node IDs, so adjacent
triangles use exactly the same crossing.  The resulting line segments form a
$C^0$, piecewise-linear interface $\Gamma_h$.  Facet normals may jump at
vertices, as expected for a P1 reconstruction; continuity of the interface
position, rather than continuity of its normal, is the essential requirement.

Segments shorter than
\texttt{weakInterfaceTraceMinSegmentFraction} times the minimum in-plane cell
width are discarded.  Unlike the previous cut-cell prototype, they are not
replaced by independent nodal ties.

\subsection{P1 vertex mortar}
\index{mortar method!weak discontinuity projection}
\index{weak discontinuity projection!local projection}

A scalar P1 multiplier hat $\lambda_a$ is associated with every reconstructed
interface vertex $a$.  It is nonzero only on the segments incident to that
vertex.  The coefficient coupling the multiplier to grid node $I$ is
\begin{equation}
  c_{aI} =
  \sum_{e\in\mathcal{E}(a)}\int_e \lambda_a(\mathbf{x})N_I(\mathbf{x})\,\mathrm{d}s.
  \label{eq:weak-trace-mortar-coefficient}
\end{equation}
The integral is evaluated with one- or two-point Gauss quadrature on each
segment; two points are the default.

For \texttt{SinglePointBSpline}, $N_I$ in
Eq.~\eqref{eq:weak-trace-mortar-coefficient} is the full tensor-product cubic
B-spline trace, including all $4\times4\times4$ support nodes used by the
particle calculation.  For linear CPDI, $N_I$ is the underlying trilinear Q1
grid field.  CPDI particle-domain transfer remains otherwise unchanged.  In
particular, this prototype does not clip a CPDI domain at $\Gamma_h$ or modify
its volume.

A coefficient is retained only when both contact-group fields are active at
that node and both nodal masses exceed
\begin{equation}
  m_{\mathrm{cut}}^s =
  \max\!\left(m_{\min},
  \texttt{weakInterfaceTraceMassCutoffRatio}\,\max_I m_I^s\right),
  \qquad s\in\{A,B\}.
  \label{eq:weak-trace-mass-cutoff}
\end{equation}
The same coefficient $c_{aI}$ is used on both sides and is \emph{not}
renormalized independently.  This preserves the common physical trace being
tied.  A constraint is accepted only when the retained coefficient sum divided
by the unfiltered sum is at least
\texttt{weakInterfaceTraceMinWeight}.

For velocity component $\alpha$, the integrated trace residual and local
inverse effective mass are
\begin{align}
  g_{a\alpha} & =
  \sum_I c_{aI}\left(v_{I\alpha}^A-v_{I\alpha}^B\right),
  \label{eq:weak-trace-mortar-residual}\\
  D_a & =
  \sum_I c_{aI}^2\left(\frac{1}{m_I^A}+\frac{1}{m_I^B}\right)
  + \epsilon_D D_a^{\mathrm{ref}},
  \label{eq:weak-trace-mortar-denominator}
\end{align}
where $\epsilon_D$ is
\texttt{weakInterfaceTraceLinearSolverRegularization}.  One relaxed local
correction is
\begin{equation}
  p_{a\alpha}=-\omega\frac{g_{a\alpha}}{D_a},
  \qquad
  \omega=\texttt{weakInterfaceTraceProjectionScale},
  \label{eq:weak-trace-mortar-impulse}
\end{equation}
followed by the equal-and-opposite velocity increments
\begin{equation}
  v_{I\alpha}^A \leftarrow v_{I\alpha}^A+
    \frac{c_{aI}p_{a\alpha}}{m_I^A},
  \qquad
  v_{I\alpha}^B \leftarrow v_{I\alpha}^B-
    \frac{c_{aI}p_{a\alpha}}{m_I^B}.
  \label{eq:weak-trace-mortar-scatter}
\end{equation}
Every local update conserves linear momentum to roundoff.  Because adjacent
vertex hats overlap, the constraints are visited in alternating forward and
reverse order for the requested number of sweeps.  This is a projected
Gauss--Seidel interface iteration composed entirely of local corrections.  It
does not assemble or factor a global matrix.  The iteration can stop early when
the relative RMS mortar residual falls below
\texttt{weakInterfaceTraceLinearSolverTolerance}.

After projection, the net velocity change is converted to
\texttt{gridWeakInterfaceTraceForce}$=m(\mathbf{v}^{\mathrm{proj}}-
\mathbf{v}^{\mathrm{trial}})/\Delta t$ and applied through the standard nodal
momentum, acceleration, and velocity-increment fields.  This preserves the
usual explicit-step bookkeeping and makes the interface impulse available as a
diagnostic force.

\subsection{Placement in the explicit step and contact interaction}
\index{weak discontinuity projection!solver step}
\index{contact!weak discontinuity projection}

The projection is applied after ordinary grid dynamics and contact and before
prescribed deformation and essential boundary conditions.  When
\texttt{weakInterfaceTraceSuppressNodalContact=1}, ordinary nodal contact is
skipped for selected contact-group pairs.  Otherwise, nodal no-slip contact can
collapse the two fields toward the single-field mixed-cell kinematics that the
mortar treatment is intended to avoid.

The suppression therefore assumes that the reconstructed mortar covers the
entire bonded interface.  Missing surface samples, an interface too close to a
B-spline support boundary, insufficient common nodal mass, or rejected short
segments can leave part of a selected pair uncoupled.  The reconstruction and
constraint-count log message and the diagnostic fields listed below should be
checked in every new problem.

\subsection{Input controls and diagnostic fields}
\label{subsec:weak-trace-controls}
\index{weak discontinuity projection!input controls}

\begingroup\small
\begin{longtable}{>{\raggedright\arraybackslash}p{0.40\linewidth}>{\raggedright\arraybackslash}p{0.52\linewidth}}
\caption{Experimental $C^0$ vertex-mortar controls.}\label{tab:weak-trace-controls}\\
\toprule
Attribute & Effect \\
\midrule
\endfirsthead
\toprule
Attribute & Effect \\
\midrule
\endhead
\texttt{enableWeakInterfaceTraceProjection} & Enables the prescribed-surface weak-discontinuity projection.  Disabled by default. \\
\texttt{weakInterfaceTracePairs} & Optional two-column list of contact-group pairs.  Empty means all distinct group pairs. \\
\texttt{weakInterfaceTraceSurfaceFitNeighbors} & Number of nearest oriented surface samples used in Eq.~\eqref{eq:weak-trace-nodal-phi}.  Default: 4. \\
\texttt{weakInterfaceTraceMortarQuadratureOrder} & Segment Gauss rule; supported values are 1 and 2.  Default: 2. \\
\texttt{weakInterfaceTraceMassCutoffRatio} & Relative side-mass cutoff in Eq.~\eqref{eq:weak-trace-mass-cutoff}.  Default: $10^{-12}$. \\
\texttt{weakInterfaceTraceMinWeight} & Minimum retained common-support fraction for accepting a vertex constraint. \\
\texttt{weakInterfaceTraceCutTolerance} & Nodal $\phi$ zero tolerance relative to the in-plane cell diagonal. \\
\texttt{weakInterfaceTraceMinSegmentFraction} & Minimum P1 segment length relative to the minimum in-plane cell width. \\
\texttt{weakInterfaceTraceProjectionScale} & Local relaxation factor $\omega$ in Eq.~\eqref{eq:weak-trace-mortar-impulse}.  Default: 0.6. \\
\texttt{weakInterfaceTraceProjectionIterations} & Maximum number of alternating local projection sweeps.  Default: 30. \\
\texttt{weakInterfaceTraceLinearSolverTolerance} & Relative RMS mortar-residual tolerance for early termination. \\
\texttt{weakInterfaceTraceLinearSolverRegularization} & Dimensionless local denominator regularization $\epsilon_D$. \\
\texttt{weakInterfaceTraceSuppressNodalContact} & Suppresses ordinary nodal contact for mortar-owned group pairs when nonzero. \\
\texttt{weakInterfaceTraceGapStabilization} & Reserved; must currently be zero. \\
\texttt{weakInterfaceTraceMaxGapCorrectionVelocity} & Reserved; must currently be zero. \\
\bottomrule
\end{longtable}
\endgroup

The principal transient diagnostics are
\texttt{gridWeakInterfaceTraceActive},
\texttt{gridWeakInterfaceTraceForce},
\texttt{gridWeakInterfaceTracePoint},
\texttt{gridWeakInterfaceTraceSupportWeight},
\texttt{gridWeakInterfaceTraceVelocityJump}, and
\texttt{gridWeakInterfaceTraceVelocityJumpPost}.  The solver log additionally
reports the numbers of surface markers, reconstructed vertices and segments,
accepted constraints, rejection reasons, performed sweeps, residual reduction,
kinetic-energy change, and momentum change.

\subsection{Conservation and EOS state}
\index{equation of state!weak discontinuity projection}
\index{mass conservation!weak discontinuity projection}

The mortar algorithm changes only nodal velocities through balanced impulses.
It does not rescale particle mass or current/reference volume, does not alter
particle density, and does not apply a cut-cell zeroth-moment volume correction.
Consequently, it introduces no artificial density change into an EOS call.  A
higher density of ordinary particles near the interface may be generated in PFW
as a separate discretization refinement if surface sampling or CPDI-domain
resolution proves insufficient; such refinement is not part of the current
solver algorithm.

\subsection{Current limitations and verification guidance}
\index{weak discontinuity projection!limitations}

The current implementation has the following deliberate restrictions:
serial execution, CPU execution, plane strain, prescribed
\texttt{WeakDiscontinuity} surface samples, contact-group field splitting, and
either all cubic B-spline particles or all linear CPDI particles.  Linear
single-point MPM is rejected because the mortar tie does not cure its underlying
cell-crossing instability.  Mixed B-spline/CPDI subregions, MPI reconstruction,
periodic-interface stitching, three-dimensional surface reconstruction,
automatic material-interface discovery, gap stabilization, debonding, and CPDI
domain clipping remain future work.

Verification should examine local stress error and false inelasticity in a band
around the interface, not only the far-field load--displacement response.  For a
curved interface, the desired qualitative behavior is a smooth error associated
with the P1 geometric approximation, without the alternating grid-scale halo
created by incompatible independent cell-gradient corrections.  The pre- and
post-projection mortar velocity jumps, momentum change, and interface
reconstruction counts should be reported together with local
$\sigma_{nn}$, tangential stress, von Mises stress, and equivalent plastic
strain.
